%% -*- coding:utf-8 -*-
\documentclass[fleqn,10pt,a4paper]{article}


\usepackage{times}
\usepackage{amssymb}

%\usepackage[T1]{fontenc}
%\usepackage[utf8x]{inputenc}

\usepackage[indentfirst=false]{xeCJK}
\setCJKmainfont{SimSun}

\usepackage{german,article-ex,makros.2e,de-date,my-gb4e-article,float,jambox}


\usepackage[sectionbib]{natbib}

% \usepackage[bookmarks=true,hyperindex=true,breaklinks=true,draft=false,plainpages=false,hyperfootnotes=false,%
% pdfauthor={Stefan Müller},%
% pdftitle={Logik für Linguisten},%
% pdfkeywords={Logik}
% ,ps2pdf=true  %ohne diesen Treiber geht der Zeilenumbruch in URLs
% ]{hyperref}% for pdf files


\usepackage{mycommands}

% war 2018
\title{Klausur: BA Syntax 2024 (1)}
\author{Prof.\ Dr.\,Stefan Müller}


\begin{document}
\maketitle


\noindent
Name und Vorname:

\noindent
Matrikelnummer: 

\medskip
\noindent
Hinweise:
\begin{itemize}
%\item Bitte alle Lösungen direkt auf den Aufgabenblättern notieren.
\item Hilfsmittel sind nicht zugelassen.
\item Für die Bearbeitung der Klausur stehen 36 Minuten zur Verfügung.
\item Bitte schreiben Sie Ihren Namen und die Matrikelnummer auf alle Blätter.
\item Insgesamt sind 40 Punkte erreichbar. 
%% Die Punkteverteilung gibt Ihnen also einen Anhaltspunkt,
%%   wieviel Zeit Sie auf eine Aufgabe verwenden sollten. Jeder Teilnehmer hat zu Beginn 3 Punkte. Bei
%%   Multiple-Choice-Aufgaben wird pro falscher Antwort ein Punkt abgezogen.

%\item Zum Bestehen der Klausur genügt die Hälfte (50) der maximal erreichbaren Punkte.
\end{itemize}


%% \begin{enumerate}
%% \item Welche der angeführten Sprachen ist Isländisch? [4]
%% \begin{itemize}
%% \item[$\bigcirc$] at Peter har forsøgt at reparere bilen
%% \item[$\bigcirc$] Ikh meyn  az   haynt hot Max geleyent dos bukh.
%% \item[$\bigcirc$] Oft var   talað      um   þennan mann.
%% \item[$\bigcirc$] 对于一门发展中的、涉及内容广泛的学科而言这是正常的，但长期下去，又会对学科的发展产生不利影响。 
%% \end{itemize}

%% \item Wenn einen Aussagesatz mit der Abfolge XP V YP V aus einer beliebigen germanischen Sprache vorgelegt
%%   bekommen, was können Sie dann über die typologische Einordnung sagen? [4]
%% \begin{itemize}
%% \item[$\bigcirc$] Es handelt sich um eine V2-Sprache.
%% \item[$\bigcirc$] Es handelt sich um eine SVO-Sprache.
%% \item[$\bigcirc$] Es handelt sich um eine SOV-Sprache.
%% \item[$\bigcirc$] Man kann die Sprache nicht einordnen.
%% \end{itemize}

%% \newpage
%% \item Welcher Test funktioniert für den Subjektstatus in allen germanischen Sprachen? [4]
%% \begin{itemize}
%% \item[$\bigcirc$] Kontrollierbarkeit
%% \item[$\bigcirc$] Subjekt-Verb-Kongruenz
%% \item[$\bigcirc$] Stellung nach dem finiten Verb, wenn eine andere Konstituente vorangestellt wurde
%% \item[$\bigcirc$] Es gibt keinen Test, der in allen Sprachen funktioniert.
%% \end{itemize}

%% \item Die richtige Struktur für einen dänischen eingebetteten Satz ist. [4]
%% \begin{itemize}
%% \item[$\bigcirc$] 
%% \end{itemize}

%% \item Welches Phänomen tritt nur bei den germanischen OV-Sprachen und nicht bei den anderen
%%   germanischen Sprachen auf? [4]
%% \begin{itemize}
%% \item[$\bigcirc$] V2
%% \item[$\bigcirc$] Scrambling
%% \item[$\bigcirc$] Expletive Subjekte
%% \item[$\bigcirc$] keines der drei genannten
%% \end{itemize}

%% %\item Wobei handelt es sich bei dem \emph{es} in (\mex{1})?
%% %\ea

%% \item Welche Kasus kommen im Deutschen als lexikalische Kasus vor? [4]
%% \begin{itemize}
%% \item[$\bigcirc$] alle vier Kasus
%% \item[$\bigcirc$] nur der Dativ
%% \item[$\bigcirc$] nur der Genitiv
%% \item[$\bigcirc$] nur der Akkusativ
%% \end{itemize}

%% \item Welche Kasus kommen im Isländischen als Subjektskasus vor? [4]
%% \begin{itemize}
%% \item[$\bigcirc$] alle vier Kasus
%% \item[$\bigcirc$] nur der Nominativ
%% \item[$\bigcirc$] nur der Nominativ und der Dativ
%% \item[$\bigcirc$] nur der Nominativ, der Genitiv und der Dativ
%% \end{itemize}

% 6*4 = 24 + 6 = 30




%% \item Bestimmen Sie die Valenz der folgenden Wörter (Mehrfachnennungen möglich): [10]
%% \eal
%% \ex Baum
%% \ex lachen
%% \ex rasieren
%% \ex überlassen
%% \ex Bild
%% \zl

%% \item Analysieren Sie die folgenden drei Sätze (Angabe der Baumstruktur mit Kategorien und Valenzmerkmalen): [15]
%% \eal
%% \ex dass er es liest
%% \ex 
%% \gll at han læser bogen\\
%%      dass er liest Buch.{\sc def}\\
%% \ex that he reads it
%% \zl

\bigskip

\begin{enumerate}
\item Nennen Sie mindestens 10 germanische Sprachen. (5 Punkte)
% Für je zwei einen Punkt:
% Englisch, Deutsch, Niederländisch, Dänisch, Schwedisch, Norwegisch, Jiddisch, Afrikans,
% Niederländisch, Isländisch
\item Analysieren Sie die folgenden Sätze (Angabe der Baumstruktur mit Kategorien und
  Valenzmerkmalen). NPen und PPen können Sie abkürzen. (24 Punkte)
\eal
\ex dass Aicke auf das Buch wartet
%5
\ex that Aicke waits for the book
%5
\ex Aicke wartet auf das Buch.
%10
\ex 
\gll Aicke venter på bogen.\\
     Aicke wartet auf Buch.\textsc{def}\\\jambox{(Danish)}
%10
\zl
\item Erklären Sie die Analyse des Passivs im Deutschen im in der Vorlesung besprochenen
  theoretischen Rahmen. (11 Punkte)
\end{enumerate}




%% \section*{Zusatzaufgaben}

%% \begin{enumerate}
%% \item Auf welchen Kontinenten werden germanische Sprachen (außer Englisch) als Amtssprachen gesprochen? [4] 
%% \begin{itemize}
%% \item[$\bigcirc$] nur in Europa
%% \item[$\bigcirc$] in Europa, Südamerika und Nordamerika
%% \item[$\bigcirc$] in Europa, Afrika, Südamerika und Nordamerika
%% \item[$\bigcirc$] in Europa, Afrika, Südamerika, Nordamerika und Asien
%% % Mit Englisch ist es Asien auch Singapur, Hongkong, Indien 
%% \end{itemize}

%% \item Skizzieren Sie die Analyse des Passivs im Deutschen oder im Isländischen. [10]
%% \item Nennen Sie mindestens zwei germanische Sprachen, die einen Verbalkomplex bilden. Wie kann man
%%   diese Verbalkomplexe analysieren? [10]

%% \item Wie kann man erklären, dass es im Dänischen die beiden Abfolgen in (\mex{1}) gibt? [5]
%% \eal
%% \ex
%% \gll  Jens læser ikke bogen.\\
%%       Jens liest nicht  Buch.{\sc def}\\
%% \ex
%% \gll at Jens ikke læser bogen\\
%%      dass Jens nicht liest Buch.{\sc def}\\
%% \zl

%% \end{enumerate}





\end{document}





